% HD-N24. GPT-6 Astra Pro /root/imprimitive, 21 November 2026.
% Inherited effort; numerical setting not exposed.
\section{HD-N24: every partial normalization of the positive graph union}
\label{sec:partial-normalization-graph-union}

Retain any of the nodal arrangements $Z\subset A$ with the
actual directions and multiplicities of HD-N23. Write $Q$ for
its node set, $q=2350080$, and $N=9408$ for its number of
connected abelian-threefold components. For every subset
$S\subset Q$, form the actual partial normalization
$\nu_S:Z_S\to Z$ which separates exactly the nodes in $S$.
The perfect complex
\[
 F_S=\left[\mathcal O_A\longrightarrow
                    \nu_{S,*}\mathcal O_{Z_S}\right]
                  \quad\hbox{in degrees }0,1                \tag{N24.1}
\]
uses the literal restriction map. It is gluable and simple, and
its whole character is
\[
 \operatorname{ch}(F_S)
 =1-32\theta^3-120T+\frac{q-|S|}{5040}\theta^6,
 \qquad\chi(F_S,F_S)=-2|S|.                                \tag{N24.2}
\]
For every actual coherent action of a faithful finite point
subgroup on $F_S$, its invariant degree-two categorical trace
has a kernel of dimension at least $79$. This conclusion also
holds after actual isogeny pullback, line tensor, and
Fourier--Mukai equivalences, with their corresponding actual
point actions. It includes $S=\varnothing$, the original ideal,
and $S=Q$, the full normalization restriction complex. It
concerns these actual restriction complexes; arbitrary new
extension classes are not included.

\paragraph{1. The finite algebra and actual cone.}
Let $C_\nu$ denote the translated connected abelian components,
and let $j_\nu:C_\nu\hookrightarrow A$. Define the finite
$\mathcal O_Z$-algebra
\[
 \mathcal B_S=
 \ker\left(\bigoplus_\nu j_{\nu,*}\mathcal O_{C_\nu}
        \longrightarrow\bigoplus_{p\in Q\setminus S}\mathcal O_p
                                                \right),    \tag{N24.3}
\]
where each component of the map is the difference of the two
restriction values at that node. The kernel is a subalgebra:
the two branch values agree after sums and products. Put
$Z_S=\operatorname{Spec}_Z\mathcal B_S$. It is the stated
partial normalization, since at a node in $S$ its local algebra
is the direct sum of the two smooth branch rings, and at every
other node it is their fiber product over the residue field.
There is an exact sequence of actual sheaves
\[
 0\longrightarrow\mathcal O_Z\longrightarrow\mathcal B_S
             \longrightarrow\bigoplus_{p\in S}\mathcal O_p
                         \longrightarrow0.                 \tag{N24.4}
\]
Consequently
\[
 \mathcal H^0(F_S)=I_Z,\qquad
 \mathcal H^1(F_S)=\bigoplus_{p\in S}\mathcal O_p,\qquad
             \mathcal H^j(F_S)=0\ (j\ne0,1).               \tag{N24.5}
\]
This is also its actual Postnikov triangle
$I_Z\to F_S\to(\bigoplus_{S}\mathcal O_p)[-1]\to I_Z[1]$.
The maps in (N24.3)--(N24.4), rather than the $K$-class alone,
define the extension.

The exact sequence and the GRR calculation for the component
embeddings in HD-N23 give (N24.2). In particular
$[Z]^2=2q[\mathrm{pt}]$, and duality changes the sign of the
codimension-three term; thus the top term of
$\operatorname{ch}(F_S^\vee)\operatorname{ch}(F_S)$ is
$-2|S|[\mathrm{pt}]$. Negative self-Ext vanishes: for an
object with cohomology only in degrees zero and one, the only
possible negative group is controlled by
$\operatorname{Hom}(\mathcal H^1,\mathcal H^0)$ through the
Postnikov triangle. Here that Hom is zero because $I_Z$ is
torsion free and the other sheaf has finite support. The same
argument holds locally. Perfection also follows directly from
the finite free local complexes below and from the ordinary
ideal resolutions at the other points.

\paragraph{2. The normalized-node model and its exact free contraction.}
At a separated node choose the completed regular local ring
\[
 R=\mathbf C[[x_1,x_2,x_3,y_1,y_2,y_3]],\qquad
 U=R/(x_1,x_2,x_3),\quad V=R/(y_1,y_2,y_3),\quad Q_0=U\oplus V.
                                                               \tag{N24.6}
\]
The local object is the actual fiber of the restriction map,
\[
             F\longrightarrow R\xrightarrow{f}Q_0
                    \longrightarrow F[1],\qquad
             f(r)=(r\bmod x,r\bmod y).                    \tag{N24.7}
\]
For a triple $z=(z_1,z_2,z_3)$, let
\[
 d_1(z)=(z_1,z_2,z_3),\quad
 d_2(z)=\begin{pmatrix}-z_2&-z_3&0\\z_1&0&-z_3\\0&z_1&z_2\end{pmatrix},
 \quad d_3(z)=(z_3,-z_2,z_1)^t                              \tag{N24.8}
\]
be the Koszul differentials. The cone of the two Koszul
resolutions of $Q_0$, shifted by $[-1]$, has free ranks
$2,6,7,2$ in degrees $-2,-1,0,1$; its last map is
\[
                  (r,u,v)\longmapsto
                   (r-d_1(x)u,\ r-d_1(y)v).                \tag{N24.9}
\]
It contracts to the free complex of ranks $2,6,6,1$ in those
degrees, with differentials
\[
 -\operatorname{diag}(d_3(x),d_3(y)),\quad
 -\operatorname{diag}(d_2(x),d_2(y)),\quad
                         (d_1(x),-d_1(y)).                 \tag{N24.10}
\]
The projection is the identity in negative degrees,
$(r,u,v)\mapsto(u,v)$ in degree zero, and
$(a,b)\mapsto b-a$ in degree one. Its inclusion is the identity
in negative degrees, $(u,v)\mapsto(d_1(x)u,u,v)$ in degree
zero, and $b\mapsto(0,b)$ in degree one. Their composite on
the smaller complex is the identity. On the larger complex,
identity minus the opposite composite is $dh+hd$, where the
only nonzero homotopy is
$h(a,b)=(a,0,0)$ from degree one to degree zero.
These are literal chain maps over $R$ and prove the claimed
quasi-isomorphism.

\paragraph{3. The actual Hom boundary maps.}
The Koszul resolutions give
\[
 \operatorname{Ext}_R^j(Q_0,R)=
 \begin{cases}Q_0,&j=3,\\0,&j\ne3,\end{cases}
 \qquad
 \operatorname{Hom}_R(Q_0,Q_0)=Q_0,
 \qquad \operatorname{Hom}_R(R,Q_0)=Q_0.                    \tag{N24.11}
\]
The last two modules are identified by precomposition with
$f$: a pair of branch multiplication maps restricts to that
same pair. In particular this map is the identity, not a
vanishing boundary presumed from dimensions. Applying
$R\operatorname{Hom}(-,R)$ to (N24.7) gives
\[
 \operatorname{Hom}(F,R)=R,\qquad
 \operatorname{Ext}^1(F,R)=0,\qquad
 \operatorname{Ext}^2(F,R)=Q_0.                             \tag{N24.12}
\]
Applying $R\operatorname{Hom}(-,Q_0)$ and using the identity
map in (N24.11) gives canonical connecting isomorphisms
\[
 \operatorname{Hom}(F,Q_0)=\operatorname{Ext}^1(Q_0,Q_0),
 \qquad
 \operatorname{Ext}^j(F,Q_0)=\operatorname{Ext}^{j+1}(Q_0,Q_0)
                       \quad(j\ge1),                       \tag{N24.13}
\]
and $\operatorname{Hom}(F,Q_0[-1])=0$.
Every map $F\to R$ is a scalar multiple of the first triangle
map, by (N24.12). Its composite with $f$ is therefore zero.
The long exact sequence from $R\operatorname{Hom}(F,-)$ now
proves, by the displayed actual maps,
\[
 \operatorname{Hom}(F,F)=R,\qquad
 \operatorname{Ext}^1(F,F)=\operatorname{Ext}^1(Q_0,Q_0)
                                  =U^3\oplus V^3.          \tag{N24.14}
\]
Here the cross groups $\operatorname{Ext}^1(U,V)$ and its
reverse vanish: applying the $x$-Koszul resolution to $V$
is the regular three-variable Koszul cochain complex, whose
only cohomology is $\mathbf C$ in degree three.
The diagonal degree-one groups are the usual conormal-dual
modules supplied by the same Koszul resolution.

For precision, the punctual degree-two term of the original
ideal really disappears. The next exact sequence is
\[
 0\longrightarrow\operatorname{Ext}^2(Q_0,Q_0)
 \longrightarrow\operatorname{Ext}^2(F,F)\longrightarrow Q_0
 \xrightarrow{f_*}\operatorname{Ext}^3(Q_0,Q_0).            \tag{N24.15}
\]
The projection of $f_*$ onto the two diagonal top Koszul
Ext modules is the identity of $Q_0$, so $f_*$ is injective.
Thus
\[
             \operatorname{Ext}^2(F,F)=U^3\oplus V^3.       \tag{N24.16}
\]
There is no punctual nine-dimensional summand here. This local
improvement does not remove the different global term below.
At a node not in $S$, $F_S=I_Z$ and HD-N23 gives the same
$\mathcal E xt^1$ module $U^3\oplus V^3$, while retaining
its separately computed degree-two punctual term.

\paragraph{4. Canonical gluing of the first Ext sheaf.}
Let $j:A\setminus Q\hookrightarrow A$. On this open set,
$F_S$ is the actual ideal of the disjoint smooth components.
Its first self-Ext sheaf is canonically their normal bundle,
by the ideal exact sequence and the conormal-dual map. Set
\[
 \mathcal E=\mathcal E xt^1(F_S,F_S),\qquad
                   \mathcal M=\bigoplus_\nu j_{\nu,*}N_{C_\nu/A}.
                                                               \tag{N24.17}
\]
Their canonical identifications off $Q$ give
$j^*\mathcal E\simeq j^*\mathcal M$. At every node, the
completed stalk of each is $U^3\oplus V^3$, by (N24.14)
and HD-N23. Both stalks have depth at least three. One can
see this without asserting invariance of an unspecified
splitting: the parameters $x_i+y_i$, $1\le i\le3$, form
a regular sequence on both branch modules. They may be
chosen in the actual local ring by taking generators of the
two transverse branch ideals. Faithfully flat completion
then checks the same nonzerodivisor assertions before
completion.

The exact local-cohomology restriction sequence and depth
vanishing therefore give canonical isomorphisms
\[
 \mathcal E\xrightarrow{\sim}j_*j^*\mathcal E,
 \qquad\mathcal M\xrightarrow{\sim}j_*j^*\mathcal M.       \tag{N24.18}
\]
The precise primary locators are Stacks Project,
Lemma47.11.1, \href{https://stacks.math.columbia.edu/tag/0AVZ}{Tag0AVZ},
and Lemma51.2.2,
\href{https://stacks.math.columbia.edu/tag/0DWR}{Tag0DWR};
faithful completion is Lemma10.97.3, Tag00MC.
Equation(N24.18) extends the smooth normal map uniquely and
proves the actual global identity
\[
 \mathcal E xt^1(F_S,F_S)\simeq
                    \bigoplus_\nu j_{\nu,*}N_{C_\nu/A}.
                                                               \tag{N24.19}
\]
It is natural for every translation preserving the arrangement
and $S$. No independently chosen node splittings are glued.
Each normal bundle is the actual trivial rank-three quotient
bundle of the abelian embedding. Consequently
\[
             h^1(A,\mathcal E xt^1(F_S,F_S))=9N.            \tag{N24.20}
\]
The canonical scalar map
$\mathcal O_A\to\mathcal H om(F_S,F_S)$ is an isomorphism
by (N24.14), the ideal calculation at the other nodes, and
the smooth-point calculation. Local negative self-Ext vanishes
as above. Hence global degree-zero self-Ext equals
$H^0(A,\mathcal O_A)=\mathbf C$, proving simplicity.

\paragraph{5. Survival and every actual coherent point action.}
The local-to-global self-Ext spectral sequence is first quadrant
and converges, since the complexes are perfect and their
negative Ext sheaves vanish. From the position $(1,1)$, the
only outgoing differential is
\[
 H^1(\mathcal E xt^1(F_S,F_S))\longrightarrow
                         H^3(\mathcal O_A),                \tag{N24.21}
\]
with target dimension $20$; there are no incoming differentials.
Thus, without a degeneration assumption,
\[
 \dim\operatorname{Ext}^2(F_S,F_S)\ge9N-20,
 \qquad\dim\ker\sigma^2_{F_S}\ge9N-515,                    \tag{N24.22}
\]
where the full trace target has dimension $495$ by HD-N23.
For the arrangement upstairs the latter lower bound is $84157$.

If $T_a^*F_S\otimes P\simeq F_S$ for $P\in\widehat A$,
taking $\mathcal H^0$ and its double dual gives
$P\simeq\mathcal O_A$. The induced ideal isomorphism extends
to a scalar automorphism of $\mathcal O_A$, so its natural
ideal inclusions have the same zero scheme. Thus $a$ preserves
$Z$. Taking $\mathcal H^1$ also makes it preserve the selected
node set $S$. Conversely, a translation preserving $(Z,S)$
preserves the literal algebra (N24.3) and restriction map,
so gives its canonical action on $F_S$.
The full point stabilizer is therefore the translation group
$H_S=\operatorname{Stab}(Z,S)$, a subgroup of the actual
$H$ in HD-N23. In particular $|H_S|\le144$.

For a faithful coherent point subgroup $G\subset H_S$, take
the actual finite etale quotient $p:A\to A'=A/G$.
The component images are smooth translated abelian threefolds;
their number $N'$ is the number of component orbits, so
\[
                      N'\ge\left\lceil\frac{9408}{|G|}\right\rceil
                               \ge66.                      \tag{N24.23}
\]
Two different components of one orbit are parallel and disjoint.
Thus no pair of branches of a node is identified into one
component, and $Z'\subset A'$ again has only transverse double
points and no triples. The invariant subset $S$ descends to
its node subset $S'$. Finite etale pullback of the actual kernel
algebra (N24.3) identifies the partial-normalization complex
$F'_{S'}$ downstairs with $F_S$ upstairs:
\[
                              p^*F'_{S'}=F_S.               \tag{N24.24}
\]
This also follows directly by descent on its two-branch local
algebras; the kernel formula fixes the restriction map.
Apply (N24.22) to that actual complex on $A'$. Its trace
kernel has dimension at least
\[
                               9\cdot66-515=79.             \tag{N24.25}
\]
Pullback on Ext is injective, split by the adjunction trace
divided by $\deg p$, and the Atiyah class and categorical
trace commute with etale pullback. The resulting classes
are invariant under the canonical deck action. Simplicity
proved above implies that every other choice of the coherent
isomorphisms on the same point subgroup differs by scalars;
they cancel on Ext conjugation. The same classes are therefore
invariant for the given coherent action. HD-N12's identity-sector
comparison identifies their ambient trace kernel with the weak
equivariant trace kernel, including the stated scalar coherence.

Finally let $f:B\to A$ be an actual isogeny. The same
cohomology and double-dual argument identifies the full point
stabilizer of $f^*F_S$ with $f^{-1}(H_S)$. Apply the preceding
quotient construction to $A/H_S$ and pull its trace-kernel
classes through $B\to A\to A/H_S$. The actual unit/trace
splitting keeps them nonzero, and they are invariant under
$f^{-1}(H_S)$ and hence every subgroup. The pullback complex
has the same local models and scalar endomorphisms, so all
coherent choices give the same conjugation on Ext. Actual
line-tensor and Fourier--Mukai equivalences preserve the
trace and conjugate point actions by the maps of HD-N12 and
HD-N17. These exact maps prove the full stated scope.

\paragraph{Verification.}
The certificate is \path{partial_normalization_graph_certificate.py}.
It checks the displayed Koszul chain maps, the explicit cone
projection/inclusion/homotopy identities, and all final numerical
bounds. The Hom boundary maps, canonical depth extension,
local-to-global survival and actual equivariant descent are
proved above; finite computations do not replace them.
